\section{Overview}
\label{m:sec:overview}

This chapter introduces the mathematical framework and the methods used
throughout the thesis. Its concern is the machinery rather than any single
application: the process classes in which the models live, and the stock of
exact arguments by which those processes are analysed. Later chapters build
their models from the objects defined here, and appeal to the methods developed
here; the present one aims to make both available in a form that can be cited
rather than re-derived.

Two questions organise everything below, and it is worth naming them because
they are not the same question. The first is a \emph{distributional} question:
given the process, what is the law of the population --- its mean, its variance,
its full generating function --- at a given time, or in the limit? The second is
an \emph{absorption} question: does the process reach a given state at all, and
if so with what probability, or after what time? Extinction, rupture and
establishment are all absorption questions; growth and quasi-stationarity are
distributional ones. The methods of \cref{m:sec:methods} split accordingly.
First-step analysis and its continuous-time companions answer absorption
questions; generating functions and moment equations answer distributional ones;
and where a problem mixes the two --- survival \emph{conditioned on} not being
absorbed, for instance --- the two stocks of methods are used together, as in
\cref{m:sec:condmean}.

\Cref{m:sec:markov} defines the process classes. Discrete-time Markov chains are
introduced through simple random walks and Galton--Watson processes; the latter
is the ancestor of every branching model in the thesis, and is treated here only
as far as its definition and one-step structure, with the critical theory in
\cref{m:app:critical} and the deep analysis of its quasi-stationary amplitude
in \ChA. Continuous-time chains follow, through the Poisson process and the
birth--death and birth--death--catastrophe processes, the latter carrying the
second absorbing mechanism --- abrupt rupture --- that several later chapters
need. Time-inhomogeneous chains, whose rates vary with the clock, are set up
with the logistic speciation model as the worked instance; and coupled
ODE--CTMC systems, in which a deterministic flow and a jump process drive one
another, close the section with the rupture-into-a-medium schematic on which the
compartment chapters are built.

\Cref{m:sec:methods} then develops the methods, again in two time domains. The
discrete-time stock --- generating functions, functional iteration, first-step
analysis --- is set out briefly, because its deeper development belongs to \ChA,
where the binary Galton--Watson process is analysed by Koenigs-function methods.
The continuous-time stock is the main technical weight of the chapter: backward
equations, mean population size, variance and higher moments, hitting
probabilities, extinction probabilities, and conditional means, all developed
with the birth--death process as the running example. The conditional-mean
subsection is where the two halves of the thesis meet: it defines the discrete
constant $A(p)$ and the quasi-stationary mean $1/A(p)$, derives the elementary
continuous-time constant $\Ac(p)$ in full, and hands the analysis of $A(p)$
itself --- product and series representations, bounds, near-critical asymptotic,
closed-form searches, and the dynamical obstruction to any closed form --- to
\ChA.

One method crosses the division and receives its own treatment in
\cref{m:sec:moc}: the method of characteristics, by which the first-order
partial differential equations satisfied by generating functions are solved. It
is presented as a four-step recipe and worked through on the absorption--death
model; the harder members of the family, including the absorption--birth--death
model whose closed form is an original result of this thesis, are solved in
\cref{m:app:moc} by the identical procedure, supported by the integral identity
of \cref{m:app:hyp} and the coefficient-extraction machinery of
\cref{m:app:extract}. The remaining appendices hold material displaced from the
main text by the chapter's framework-and-methods organisation --- the critical
theory of branching, the arithmetic of founding cohorts, the quasi-stationary
variance, and the discrete catastrophe calculation --- so that nothing the later
chapters need is lost.

Two conventions run throughout. Processes are specified by their rates and
nothing else: \cref{m:eq:infinitesimal} for continuous time, the offspring law
or transition matrix for discrete time, and every property is read off from
that specification by one of the methods above. And the voice is deliberately
abstract: the mathematics of a compartment that may rupture does not depend on
what the compartment contains, and the chapter is written so that later
chapters are free to supply the biology.
